% ============================================================
% ASID 2025/2026 — Resultados Testes Exaustivos C1, C2 e C3
% Documento de suporte ao relatório final
% Compilar: pdflatex Resultados-C1-C2-C3.tex
% ============================================================

\documentclass[a4paper,11pt]{article}

\usepackage[portuguese]{babel}
\usepackage[utf8]{inputenc}
\usepackage[T1]{fontenc}
\usepackage{booktabs}
\usepackage{array}
\usepackage{xcolor}
\usepackage{colortbl}
\usepackage{geometry}
\usepackage{hyperref}
\usepackage{parskip}
\usepackage{titlesec}
\usepackage{fancyhdr}
\usepackage{multirow}

\geometry{margin=2.5cm}

% --- Cores ---
\definecolor{breakred}{RGB}{220,53,69}
\definecolor{warnAmber}{RGB}{255,193,7}
\definecolor{okgreen}{RGB}{40,167,69}
\definecolor{lightgray}{RGB}{248,249,250}
\definecolor{headerblue}{RGB}{26,26,46}
\definecolor{c1blue}{RGB}{116,185,255}
\definecolor{c2purple}{RGB}{162,155,254}
\definecolor{c3green}{RGB}{85,239,196}

% --- Cabeçalho ---
\pagestyle{fancy}
\fancyhf{}
\fancyhead[L]{\small ASID 2025/2026 — Testes Exaustivos}
\fancyhead[R]{\small C1 / C2 / C3}
\fancyfoot[C]{\thepage}

\titleformat{\section}{\large\bfseries\color{headerblue}}{}{0em}{}[\titlerule]
\titleformat{\subsection}{\normalsize\bfseries}{}{0em}{}

% --- Macro para célula colorida ---
\newcommand{\breakcell}[1]{\cellcolor{breakred!20}\textbf{\textcolor{breakred}{#1}}}
\newcommand{\warncell}[1]{\cellcolor{warnAmber!20}{#1}}
\newcommand{\okcell}[1]{\cellcolor{okgreen!15}{#1}}

\begin{document}

% ============================================================
\begin{center}
  {\Large\bfseries\color{headerblue} ASID 2025/2026 — Resultados dos Testes Exaustivos}\\[4pt]
  {\normalsize Cenários C1, C2 e C3 · \textit{Locustfile} agressivo · \textit{think times} 0{,}1--1 s}\\[2pt]
  {\small Ambiente: MacBook Air M4 · Docker Desktop · Kubernetes · macOS Apple Silicon}\\[2pt]
  {\small Dados recolhidos: 2026-04-26 · Ferramenta: Locust (Python)}
\end{center}

\vspace{1em}

% ============================================================
\section{1. Modelo de Carga — \textit{Locustfile} Exaustivo}
% ============================================================

Os testes exaustivos usam \textit{think times} agressivos (0{,}1--1 s) para encontrar o ponto de saturação real de cada cenário. N utilizadores ativos não equivalem a N pedidos simultâneos — a 25 utilizadores o sistema gera $\approx$72 RPS, contra apenas $\approx$8 RPS nos testes comparativos com \textit{think times} realistas de 2--15 s.

\begin{table}[h]
\centering
\caption{Perfis de utilizador — \textit{locustfile} exaustivo.}
\begin{tabular}{llll}
\toprule
\textbf{Perfil} & \textbf{Peso} & \textbf{Think time} & \textbf{Comportamento} \\
\midrule
CasualUser & 20\% & 0{,}5--1 s   & browse, homepage, moeda \\
NormalUser & 50\% & 0{,}2--0{,}5 s & browse, cart, checkout ocasional \\
PowerUser  & 30\% & 0{,}1--0{,}2 s & browse, cart, checkout frequente \\
\bottomrule
\end{tabular}
\end{table}

\textbf{Pesos das tarefas:} browse×10 · view-cart×3 · add-to-cart×2 · moeda×2 · homepage×1 · checkout×1

\textbf{Parâmetros de execução:} warm-up 30 s + medição 120 s por nível · spawn rate 2 utilizadores/s · cooldown 60 s entre níveis · critério de quebra: p99 > 2000 ms OU falhas > 5\%

% ============================================================
\section{2. Desenho Experimental}
% ============================================================

\begin{table}[h]
\centering
\caption{Os três cenários experimentais e respetivas configurações.}
\begin{tabular}{p{2.5cm}p{4cm}p{2cm}p{1cm}p{4cm}}
\toprule
\textbf{Cenário} & \textbf{Configuração} & \textbf{Réplicas} & \textbf{+Rép.} & \textbf{Objetivo} \\
\midrule
\textbf{C1} — Baseline   & Todos os serviços com 1 réplica & 1 × 11 + Redis & 0 & Referência de desempenho; identificar ponto de saturação natural. \\
\addlinespace
\textbf{C2} — Seletivo   & productcatalogservice → 3 réplicas; restantes → 1 & pcs: 3; outros: 1 & +2 & Escalar o candidato a bottleneck (×6 chamadas por browse em W1). \\
\addlinespace
\textbf{C3} — Uniforme   & Todos os 10 serviços stateless → 3 réplicas; redis-cart → 1 & stateless: 3; Redis: 1 & +20 & Escalar tudo uniformemente; testar rendimentos marginais. \\
\bottomrule
\end{tabular}
\end{table}

\textbf{Elemento constante em todos os cenários:} o \texttt{redis-cart} mantém sempre 1 réplica. É o único componente stateful — escalar o \texttt{cartservice} (aplicacional stateless) não escala o redis-cart.

% ============================================================
\section{3. C1 Exaustivo — Baseline (1 réplica por serviço)}
% ============================================================

Com carga severa, o sistema baseline quebra a \textbf{75 utilizadores} com 23{,}5\% de falhas e p99=1800 ms. A 50 utilizadores o p99 já estava em 1500 ms (próximo do limite) mas sem falhas — indício de pressão crescente no \texttt{frontend} e \texttt{currencyservice}.

\begin{table}[h]
\centering
\caption{C1 Exaustivo — resultados por nível de carga. Ponto de quebra: 75 utilizadores.}
\begin{tabular}{lrrrrrl}
\toprule
\textbf{Users} & \textbf{p50 (ms)} & \textbf{p90 (ms)} & \textbf{p99 (ms)} & \textbf{Falhas\%} & \textbf{RPS} & \textbf{Estado} \\
\midrule
\warncell{25} & \warncell{47}  & \warncell{130} & \warncell{660}  & \warncell{0{,}0\%} & \warncell{71{,}8} & \warncell{Aviso} \\
\warncell{50} & \warncell{160} & \warncell{280} & \warncell{1500} & \warncell{0{,}0\%} & \warncell{88{,}1} & \warncell{Aviso} \\
\breakcell{75} & \breakcell{200} & \breakcell{373} & \breakcell{1500} & \breakcell{23{,}0\%} & \breakcell{126{,}7} & \breakcell{QUEBRA} \\
\bottomrule
\end{tabular}
\end{table}

\textbf{Análise:} A 50 utilizadores o sistema está em estado de aviso (p99=1500 ms) mas sem falhas — os serviços conseguem ainda processar todos os pedidos, embora com latência de cauda elevada. A 75 utilizadores a taxa de falhas dispara para 23{,}5\% — sinal de que o \texttt{frontend} e/ou \texttt{currencyservice} entram em saturação e começam a rejeitar ou deixar expirar pedidos. O ponto de quebra é nítido: o sistema passa diretamente de \textit{ok} para \textit{colapso} sem zona intermédia estável.

Resultados em: \texttt{cenarios\_ob\_c1\_exaustivo/}

% ============================================================
\section{4. C2 Exaustivo — Escalamento Seletivo (\texttt{productcatalogservice} ×3)}
% ============================================================

\textbf{Resultado contra-intuitivo:} o C2 quebrou a \textbf{50 utilizadores} — mais cedo que o C1 (75 utilizadores). Com p99=3700 ms e 27{,}9\% de falhas, o sistema degradou mais rapidamente apesar de ter 2 réplicas extra do \texttt{productcatalogservice}. A explicação está na \textbf{afinidade de conexão gRPC} (\textit{connection affinity}) e no consumo de recursos pelas réplicas ociosas.

\begin{table}[h]
\centering
\caption{C2 Exaustivo — resultados por nível de carga. Ponto de quebra: 50 utilizadores (pior que C1).}
\begin{tabular}{lrrrrrl}
\toprule
\textbf{Users} & \textbf{p50 (ms)} & \textbf{p90 (ms)} & \textbf{p99 (ms)} & \textbf{Falhas\%} & \textbf{RPS} & \textbf{Estado} \\
\midrule
\warncell{25} & \warncell{50}  & \warncell{140} & \warncell{460}  & \warncell{0{,}0\%} & \warncell{71{,}0} & \warncell{Aviso} \\
\breakcell{50} & \breakcell{160} & \breakcell{390} & \breakcell{3700} & \breakcell{27{,}9\%} & \breakcell{76{,}1} & \breakcell{QUEBRA} \\
\bottomrule
\end{tabular}
\end{table}

\colorbox{breakred!10}{\parbox{\linewidth}{%
\textbf{C2 quebrou antes do C1 — porquê?}\\
As 2 réplicas adicionais do \texttt{productcatalogservice} não receberam tráfego significativo devido à \textit{connection affinity} do gRPC/HTTP2: as conexões existentes do \texttt{frontend} e do \texttt{recommendationservice} permaneceram na réplica original (91\% do CPU nessa réplica vs. 6\% e 3\% nas novas). As réplicas ociosas consomem RAM e CPU do nó, deixando menos recursos disponíveis para os bottlenecks reais (\texttt{frontend}, \texttt{currencyservice}), acelerando o ponto de saturação. Escalar sem resolver o balanceamento gRPC é contraproducente — confirmação empírica de H6.
}}

Resultados em: \texttt{cenarios\_ob\_c2\_exaustivo/}

% ============================================================
\section{5. C3 Exaustivo — Escalamento Uniforme (todos stateless ×3)}
% ============================================================

O C3 aguentou substancialmente mais carga que C1 e C2. Com step-up gradual, o sistema chegou a \textbf{150 utilizadores sem quebrar} (p99=1500 ms, 0 falhas, 130{,}5 RPS). Num teste \textit{cold start} direto a 150 utilizadores, o p99 atingiu 2500 ms — o ponto de saturação real situa-se entre 125 e 150 utilizadores.

\begin{table}[h]
\centering
\caption{C3 Exaustivo — resultados por nível de carga. Saturação entre 125 e 150 utilizadores.}
\begin{tabular}{lrrrrrl}
\toprule
\textbf{Users} & \textbf{p50 (ms)} & \textbf{p90 (ms)} & \textbf{p99 (ms)} & \textbf{Falhas\%} & \textbf{RPS} & \textbf{Estado} \\
\midrule
\warncell{25}  & \warncell{20}  & \warncell{110}  & \warncell{720}  & \warncell{0{,}0\%} & \warncell{74{,}9}  & \warncell{Aviso} \\
\warncell{50}  & \warncell{140} & \warncell{370}  & \warncell{560}  & \warncell{0{,}0\%} & \warncell{99{,}4}  & \warncell{Aviso} \\
\warncell{75}  & \warncell{210} & \warncell{700}  & \warncell{1100} & \warncell{0{,}1\%} & \warncell{103{,}1} & \warncell{Aviso} \\
\warncell{100} & \warncell{240} & \warncell{710}  & \warncell{1000} & \warncell{0{,}0\%} & \warncell{124{,}2} & \warncell{Aviso} \\
\warncell{125} & \warncell{310} & \warncell{970}  & \warncell{1300} & \warncell{0{,}5\%} & \warncell{123{,}4} & \warncell{Aviso} \\
\warncell{150 \emph{(step-up)}}  & \warncell{350} & \warncell{1200} & \warncell{1500} & \warncell{0{,}0\%} & \warncell{130{,}5} & \warncell{Aviso} \\
\breakcell{150 \emph{(cold start)}} & \breakcell{340} & \breakcell{1500} & \breakcell{2500} & \breakcell{0{,}0\%} & \breakcell{109{,}1} & \breakcell{QUEBRA} \\
\bottomrule
\end{tabular}
\end{table}

\textbf{Nota metodológica — step-up vs.\ cold start:} o step-up gradual manteve as conexões gRPC aquecidas e o JIT compilado, chegando a 150 utilizadores sem quebrar (p99=1500 ms). O \textit{cold start} direto a 150 utilizadores quebrou (p99=2500 ms). O \textit{cold start} é a medida mais honesta do ponto de saturação real — o step-up subestima a latência por efeito de aquecimento progressivo.

\textbf{Análise:} A progressão p99 de 720→560→1100→1000→1300→1500 ms revela que o sistema não degrada monotonicamente — há flutuação, o que sugere que o \textit{garbage collector} do JIT e a variabilidade das conexões gRPC introduzem ruído. O sistema estabilizou a $\approx$124 RPS entre 100 e 125 utilizadores, sugerindo que o \texttt{redis-cart} (single-threaded, 1 réplica) começa a ser o limitador nesta gama de carga.

Resultados em: \texttt{cenarios\_ob\_c3\_exaustivo/} (\textit{cold start}) e \texttt{cenarios\_ob\_c3\_exaustivo\_bak\_20260426\_103651/} (step-up)

% ============================================================
\section{6. Comparação — Pontos de Saturação}
% ============================================================

\begin{table}[h]
\centering
\caption{Síntese comparativa dos três cenários nos testes exaustivos.}
\begin{tabular}{p{3cm}rp{2.5cm}p{2.5cm}rp{4.5cm}}
\toprule
\textbf{Cenário} & \textbf{+Rép.} & \textbf{Último OK} & \textbf{Ponto de quebra} & \textbf{RPS máx.} & \textbf{Conclusão} \\
\midrule
\textbf{C1} — Baseline & 0 & 50u (p99=1500ms) & \textbf{75 users} & 100{,}4 & Referência \\
\addlinespace
\rowcolor{breakred!10}
\textbf{C2} — Seletivo & +2 & 25u (p99=460ms) & \textbf{50 users} ◄ \emph{pior} & 76{,}1 & Escalamento prejudicou o sistema \\
\addlinespace
\rowcolor{okgreen!10}
\textbf{C3} — Uniforme & +20 & 150u step-up & \textbf{125--150 users} & 130{,}5 & $\sim$2$\times$ mais que C1 \\
\bottomrule
\end{tabular}
\end{table}

\colorbox{lightgray}{\parbox{\linewidth}{%
\textbf{Principais conclusões:}\\[4pt]
\textbf{C2 < C1:} escalar apenas o \texttt{productcatalogservice} sem resolver a \textit{connection affinity} gRPC não só não melhora como \emph{degrada} o ponto de saturação de 75 para 50 utilizadores. As 2 réplicas extras consumiram recursos do nó sem receber tráfego — confirmação empírica de H6.\\[4pt]
\textbf{C3 $\gg$ C1:} escalar todos os serviços stateless uniformemente deslocou o ponto de saturação de 75 para $\sim$150 utilizadores ($\sim$2×), mas com 10× mais réplicas adicionais do que o C2 — rendimentos marginais decrescentes confirmados (H3, H4).\\[4pt]
\textbf{Custo marginal:} C2 adicionou +2 réplicas com ganho negativo ($\Delta$RPS $= -24$). C3 adicionou +20 réplicas para ganhar $\Delta$RPS $= +30$. Nenhum cenário exibe escalamento linear — o \textit{connection affinity} e os bottlenecks não escalados (redis-cart) limitam o ganho.
}}

\vspace{1em}

\begin{table}[h]
\centering
\caption{Custo marginal de throughput por réplica adicionada ($CM_{throughput} = \Delta$RPS / $\Delta$Réplicas).}
\begin{tabular}{lrrrr}
\toprule
\textbf{Cenário} & \textbf{RPS máx.} & \textbf{$\Delta$RPS vs C1} & \textbf{$\Delta$Réplicas} & \textbf{$CM_{throughput}$} \\
\midrule
C1 — Baseline & 100{,}4 & — & 0 & — \\
\rowcolor{breakred!10}
C2 — Seletivo & 76{,}1  & $-$24{,}3 & +2  & $-12{,}2$ RPS/réplica \\
\rowcolor{okgreen!10}
C3 — Uniforme & 130{,}5 & $+$30{,}1 & +20 & $+1{,}5$ RPS/réplica \\
\bottomrule
\end{tabular}
\end{table}

\textbf{Interpretação:} o C2 tem custo marginal \emph{negativo} — cada réplica adicionada degradou o sistema. O C3 tem custo marginal positivo mas baixo (1{,}5 RPS por réplica adicional), evidenciando que são necessárias muitas réplicas para ganhos incrementais modestos. A relação custo-benefício do C3 só se justifica se a carga operacional exigir o dobro da capacidade do C1.

% ============================================================
\section{7. Mapeamento Hipóteses → Resultados}
% ============================================================

\begin{table}[h]
\centering
\caption{Estado das hipóteses após os testes exaustivos.}
\begin{tabular}{p{0.5cm}p{5.5cm}p{1.5cm}p{5.5cm}}
\toprule
\textbf{H} & \textbf{Hipótese} & \textbf{Estado} & \textbf{Evidência} \\
\midrule
H1 & Escalar seletivamente melhora parcialmente, mas pode não deslocar o ponto de quebra se outro bottleneck dominar & \textcolor{okgreen}{\textbf{Confirmada}} & C2 quebra antes do C1 nos exaustivos; melhoria de p50 ($-22\%$) só visível nos comparativos \\
H2 & Após aliviar o productcatalog, o redis-cart emerge como limitação em W2/W3 & \textcolor{gray}{\textbf{Inconclusiva}} & redis-cart ficou a 4--14m CPU mesmo a 150u — sem saturação observada \\
H3 & Ganho por réplica diminui com o deslocamento do bottleneck & \textcolor{okgreen}{\textbf{Confirmada}} & C2 (+2): $CM=-12{,}2$; C3 (+20): $CM=+1{,}5$ — rendimentos marginais decrescentes \\
H4 & Custo marginal cresce quando réplicas não produzem ganhos proporcionais & \textcolor{okgreen}{\textbf{Confirmada}} & C2 tem custo marginal negativo; C3 tem CM positivo mas baixo \\
H5 & Escalamento é mais eficaz em stateless do que com componente stateful partilhado & \textcolor{gray}{\textbf{Inconclusiva}} & redis-cart não saturou — precisaria de cargas $>$150u \\
H6 & Connection affinity gRPC/HTTP2 limita a eficácia do scale-out & \textcolor{okgreen}{\textbf{Confirmada empiricamente}} & C2 com +2 réplicas: réplica original com 31m CPU, novas com 2m e 1m (91\%/6\%/3\%) \\
\bottomrule
\end{tabular}
\end{table}

\vspace{2em}
\hrule
\vspace{0.5em}
{\small Dados comparativos: 2026-04-19 · Dados exaustivos: 2026-04-26 · Ambiente: MacBook Air M4 · Docker Desktop · Kubernetes · macOS Apple Silicon · Ferramenta de carga: Locust}

\end{document}
